\section{Partition Algebra and Countable Generation}
\label{sec:partitionalgebra}

\begin{definition}[Partition cell, partition algebra]
\label{def:partition}
Under Axiom~\ref{ax:bpsl}(iii), a \emph{cell} of the partition $\Partition_{n}$ is a maximal measurable subset $\omega\subseteq\Omega$ such that every refining partition $\Partition_{m}$ ($m\ge n$) contains either $\omega$ itself or subsets of $\omega$. The \emph{partition algebra} at depth $n$ is the Boolean algebra $\mathcal{A}_{n}$ of finite unions of cells of $\Partition_{n}$.
\end{definition}

\begin{consequence}{Countable generation}{partitionalgebra}
Under Axiom~\ref{ax:bpsl}, the Borel $\sigma$-algebra $\mathcal{B}$ of $\Omega$ is generated by a countable algebra $\mathcal{A}_{\infty}=\bigcup_{n\ge 0}\mathcal{A}_{n}$, and every observable $f\in L^{2}(\Omega,\mu)$ admits a countable partition-basis expansion.
\end{consequence}

\begin{proof}
By Axiom~\ref{ax:bpsl}(iii), $\bigvee_{n\ge 0}\Partition_{n}=\mathcal{B}$ modulo $\mu$-null sets. Each $\Partition_{n}$ contains finitely many cells by measurability and the positive-measure condition applied to $\Omega$ with $\mu(\Omega)<\infty$ (since a partition of a finite-measure set into positive-measure cells has only finitely many cells at each level). Hence $\mathcal{A}_{\infty}=\bigcup_{n}\mathcal{A}_{n}$ is a countable algebra generating $\mathcal{B}$. That every $f\in L^{2}$ admits a countable partition-basis expansion is the standard result that indicator functions of cells form a complete orthogonal system \citep{halmos1950,rudin1987}.
\end{proof}

\subsection{Concrete example of a partition algebra}

For concreteness, consider $\Omega=[0,1]^{3}$ with the uniform measure. Take $\Partition_{0}=\{\Omega\}$ (the trivial partition). Define $\Partition_{n}$ as the partition into $3^{3n}=27^{n}$ cubes of side $3^{-n}$. Each $\Partition_{n}$ refines $\Partition_{n-1}$ (each parent cube is split into 27 daughter cubes). The generated $\sigma$-algebra is the Borel algebra on $[0,1]^{3}$. This is the prototypical partition hierarchy for a three-dimensional bounded region with $b=3$ branching, to which we return in Consequence~\ref{cons:continuousembed}.

An alternative with $b=2$ branching: $\Partition_{n}$ partitions $[0,1]^{3}$ into $2^{3n}$ cubes of side $2^{-n}$. Each cube at depth $n$ is labelled by an ordered triple of $n$-digit binary fractions; the algebra is isomorphic to the Cantor algebra on ternary or binary digit sequences.

\subsection{Countable additivity}

The partition algebra $\mathcal{A}_{\infty}$ is countably additive when completed with respect to the measure $\mu$: any countable disjoint union of cells has measure equal to the sum of their measures, and any measurable set is approximable by a finite union of cells to arbitrary precision.

\section{Categorical Approximation under Finite Information}
\label{sec:categoricalapprox}

\begin{consequence}{Finite-information categorical approximation}{categorical}
Let an observer have finite information capacity $I$ bits. Then the observer can distinguish at most $2^{I}$ cells of any partition of $\Omega$; every measurement outcome is a categorical assignment to one of these cells.
\end{consequence}

\begin{proof}
An observer with $I$ bits of storage has $2^{I}$ distinct internal states; by Shannon's coding theorem \citep{shannon1948} no measurement outcome can carry more than $I$ bits of information about the system. Equivalently, the observer's reading is a function from $\Omega$ to a set of at most $2^{I}$ labels. Since the observer's reading is measurable in the partition algebra (Consequence~\ref{cons:partitionalgebra}), it factors through a partition of $\Omega$ into at most $2^{I}$ cells \citep{landauer1961}.
\end{proof}

\begin{remark}
Consequence~\ref{cons:categorical} is the axiomatic origin of what is usually called \emph{quantisation of observation}. No measurement produces a continuum of outcomes; continua exist in the unrealised phase-space background, while realised measurements populate discrete cells.
\end{remark}

\begin{corollary}[Minimum measurement energy]
\label{cor:measurementenergy}
The minimum energy required to distinguish two cells at depth $n$ is
\begin{equation}
E_{\min}(n)=\kB T_{\min}\ln 2,
\label{eq:landauer}
\end{equation}
where $T_{\min}$ is the environmental temperature and $\ln 2$ comes from binary-distinction information-theoretic cost.
\end{corollary}

\begin{proof}
Landauer's bound \citep{landauer1961}: any measurement distinguishing between two states requires logical erasure of at least one bit, which costs $\kB T\ln 2$ in energy at temperature $T$.
\end{proof}

\begin{remark}
The Landauer bound~\eqref{eq:landauer} places a lower limit on the energy cost of partition-cell identification. It also places an upper limit on the number of cells that a finite-energy observer can resolve: with energy budget $E$, the observer can perform at most $E/(\kB T\ln 2)$ bit-distinctions.
\end{remark}

\subsection{On the choice of partition base}

A recurring theme of Part I is that the branching base $b$ of the partition hierarchy is not a free parameter; it is fixed by the symmetry structure. For a system with full $SO(3)$ isotropy in three spatial dimensions, the natural branching is $b=3$ (one split per spatial axis). For a system reduced to two-dimensional isotropy, $b=2$ would be appropriate. The choice $b=3$ propagates through the Consequences: it enters the compression cost logarithm, the binding energy expression, the Fisher-metric embedding of Consequence~\ref{cons:continuousembed}, and the shell capacity formula.

The shell capacity $C(n)=2n^{2}$ does \emph{not} depend on $b$; the $2$ comes from the chirality (independent of $b$) and the $n^{2}$ comes from $\sum_{l=0}^{n-1}(2l+1)=n^{2}$. The branching $b$ enters only in the thermodynamic formulas and the Fisher embedding.

\subsection{The partition algebra as a $C^{*}$-algebra}

The indicator functions $\{\mathbf{1}_{C}\colon C\in\mathcal{A}_{\infty}\}$ form an abelian $C^{*}$-algebra under pointwise multiplication and the uniform norm. The Gelfand--Naimark theorem asserts that this algebra is isomorphic to the $C^{*}$-algebra of continuous functions on its spectrum, which in our case is the Stone space of the Boolean algebra of partition cells.

Dual to the partition algebra is the state space: each state $S\in\Omega$ corresponds to a pure state (character) of the algebra via $\chi_{S}(\mathbf{1}_{C})=1$ if $S\in C$, else $0$. Mixed states correspond to probability distributions over cells. Observations are measurements of expectation values of elements of the algebra.

This $C^{*}$-algebraic structure will not be used further in the paper, but it provides a connection to the more general framework of operator algebras in physics. In principle, all the Consequences of Parts I--V can be restated in operator-algebraic language.

\section{Partition Coordinates $(n,l,m,s)$}
\label{sec:quantumnumbers}

We now derive that the partition hierarchy of any three-dimensional bounded system admits a canonical coordinate system $(n,l,m,s)$ with specific range constraints, matching the familiar quantum numbers of atomic physics.

\begin{definition}[Principal depth]
\label{def:principal}
For each $x\in\Omega$, the \emph{principal depth} $n(x)$ is the index $n$ of the coarsest partition $\Partition_{n}$ in which $x$ occupies a cell of positive measure that is distinguished from its parent cell in $\Partition_{n-1}$. Equivalently, $n(x)=1+\max\{m: x\in\text{cell of }\Partition_{m}\text{ shared with parent in }\Partition_{m-1}\}$.
\end{definition}

\begin{lemma}[Principal depth is a positive integer]
\label{lem:nispositive}
For every $x\in\Omega$, $n(x)\in\Natural^{+}=\{1,2,3,\ldots\}$.
\end{lemma}
\begin{proof}
By Axiom~\ref{ax:bpsl}(iii), $\Partition_{0}$ is the coarsest partition; cells of $\Partition_{0}$ represent the topmost level of the hierarchy. Positive-measure cells exist at every level; thus $n(x)\ge 1$. The set is countable, giving $n(x)\in\Natural^{+}$.
\end{proof}

\subsection{Angular complexity $l$}

\begin{definition}[Angular complexity]
\label{def:angular}
At principal depth $n$, the cell containing $x$ has a geometric decomposition into an $(n-1)$-dimensional radial fibration and a two-dimensional angular profile on the $(n-1)$-sphere. The \emph{angular complexity} $l(x)$ is the index of the spherical harmonic representing the angular profile of the cell boundary.
\end{definition}

\begin{lemma}[Range constraint $l<n$]
\label{lem:langeconstraint}
At principal depth $n$, the angular complexity $l$ satisfies $0\le l\le n-1$.
\end{lemma}
\begin{proof}
The geometric decomposition of a three-dimensional bounded region at nesting depth $n$ into radial and angular components requires that the total number of independent scalar indices equals $n$. One of these is the principal depth itself; the remaining $n-1$ must be split between the radial (node count) and angular components. In a Sturm--Liouville analysis of the boundary oscillation of a bounded region \citep{arfkenweber2005}, the total nodes $n_{r}+l$ equals $n-1$, where $n_{r}\ge 0$ is the radial node count and $l\ge 0$ is the angular index. Since $n_{r}\ge 0$, we have $l\le n-1$. Combined with $l\ge 0$, this yields $0\le l\le n-1$.
\end{proof}

\subsection{Orientation $m$}

\begin{lemma}[Orientation index range]
\label{lem:mrange}
At angular complexity $l$, the irreducible representations of $SO(3)$ indexed by $l$ have dimension $2l+1$; the orientation index $m$ therefore takes values in $\{-l,-l+1,\ldots,+l\}$, a total of $2l+1$ values.
\end{lemma}
\begin{proof}
This is the standard dimension formula for the $l$th irreducible representation of $SO(3)$ \citep{wigner1959,sakurai2017}. The partition at depth $(n,l)$ inherits an $SO(3)$ action from the spherical symmetry of Axiom~\ref{ax:bpsl}(iii) on $\Omega$ taken isotropic; decomposing into $SO(3)$-irreducibles gives $2l+1$ angular sub-cells at each $(n,l)$.
\end{proof}

\begin{remark}[Zeeman observation]
The experimental verification of Lemma~\ref{lem:mrange} is the Zeeman effect \citep{zeeman1897}: placing an atom in an external magnetic field splits a spectral line of angular quantum number $l$ into exactly $2l+1$ components. The count is exact; no fractional or continuous splitting is ever observed.
\end{remark}

\subsection{Chirality $s$}

\begin{lemma}[Binary chirality]
\label{lem:chirality}
The partition at each $(n,l,m)$ admits a further binary distinction $s\in\{-\tfrac12,+\tfrac12\}$ corresponding to the two elements of the fundamental group $\pi_{1}(SO(3))=\Integer_{2}$.
\end{lemma}
\begin{proof}
$SO(3)$ is not simply connected; its fundamental group is $\Integer_{2}$ \citep{nakahara2003}. The universal cover $SU(2)\to SO(3)$ is a two-sheeted covering, so each $SO(3)$-orbit at level $(n,l,m)$ lifts to exactly two preimages in $SU(2)$. Label these preimages by $s\in\{-\tfrac12,+\tfrac12\}$ (the two half-integer weights of $SU(2)$ consistent with the binary index). This partitions the $(n,l,m)$ cell further into two; the binary chirality coordinate $s$ is well-defined at every $(n,l,m)$.
\end{proof}

\begin{consequence}{Partition coordinates}{quantumnumbers}
Every cell of the partition hierarchy at maximum depth is labelled by a quadruple $(n,l,m,s)$ with
\begin{equation}
n\in\Natural^{+},\quad l\in\{0,1,\ldots,n-1\},\quad m\in\{-l,-l+1,\ldots,+l\},\quad s\in\{-\tfrac12,+\tfrac12\}.
\label{eq:nlms}
\end{equation}
\end{consequence}

\begin{proof}
Combine Lemmas~\ref{lem:nispositive}, \ref{lem:langeconstraint}, \ref{lem:mrange}, and \ref{lem:chirality}.
\end{proof}

\begin{remark}
Consequence~\ref{cons:quantumnumbers} produces exactly the four quantum numbers of atomic spectroscopy with their standard range constraints, as observed in all atomic spectra \citep{herzberg1944,nistasd}. The derivation uses no postulate about ``orbital angular momentum'' or ``spin''; both emerge from the combined geometry and topology of $SO(3)$ acting on a bounded region.
\end{remark}

\begin{remark}[Historical note]
The ``spin'' quantum number was introduced empirically by \citet{uhlenbeckgoudsmit1925} to account for the doublet structure of atomic spectra; a theoretical derivation followed in the Dirac equation \citep{dirac1928}. Both derivations took $SO(3)$-representation structure plus the Lorentz-invariance requirement as given. In the framework of this paper, the Lorentz-invariance is itself a Consequence (Consequence~\ref{cons:lorentz}) and the doubling $s=\pm\tfrac{1}{2}$ arises from the non-triviality of $\pi_{1}(SO(3))$, which is a purely topological fact. No additional particle concept is introduced.
\end{remark}

\begin{remark}[The integer--half-integer distinction]
The integer-spin and half-integer-spin distinction---bosons and fermions in the standard terminology---is, in the deductive framework of this paper, the distinction between zero and non-zero classes in $\pi_{1}(SO(3))=\Integer_{2}$. A boson's chirality index is $s=0\pmod 1$ (trivial loop class); a fermion's is $s=\tfrac{1}{2}\pmod 1$ (non-trivial loop class). The topological origin of the distinction is manifest in the $pi_{1}$ structure.

This topological origin also explains the 720$^{\circ}$ return requirement for fermions: a fermion wave function requires two full rotations (720$^{\circ}$) to return to its original state, whereas a boson wave function returns at 360$^{\circ}$. The 720$^{\circ}$ requirement is the statement that the fermion's loop class sits in the non-trivial element of $\pi_{1}(SO(3))$, which requires a double cover (SU(2)) for continuous unwinding.
\end{remark}

\begin{proposition}[Spin--statistics structural correspondence]
\label{prop:spinstatistics}
Constituents with $s\in\{-\tfrac{1}{2},+\tfrac{1}{2}\}$ (half-integer chirality) occupy antisymmetric cells; constituents with integer chirality occupy symmetric cells.
\end{proposition}

\begin{proof}
The chirality coordinate $s$ is the element of $\pi_{1}(SO(3))=\Integer_{2}$ identifying the boundary homotopy class of the constituent. Two constituents with identical $s$ undergo combined transformation through a path of class $s+s=2s\pmod 2$. For $s=\tfrac{1}{2}$, $2s=1\pmod 2$, corresponding to a non-trivial loop in $SO(3)$: the two-constituent wave function transforms with an extra minus sign under exchange. For $s=1$ (integer), $2s=0\pmod 2$: trivial transformation, no sign. The exchange sign is anti-symmetric for half-integer and symmetric for integer chirality, reproducing the Pauli spin--statistics rule \citep{pauli1925}.
\end{proof}

\subsection{The four constraints together}

Combining Lemmas~\ref{lem:nispositive}, \ref{lem:langeconstraint}, \ref{lem:mrange}, and~\ref{lem:chirality}, the four partition coordinates are determined together. We exhibit the structure for small $n$:
\begin{center}
\begin{tabular}{c|c|c|c|c}
\toprule
$n$ & allowed $l$ & allowed $m$ per $l$ & $s$ per $(n,l,m)$ & total cells\\
\midrule
1 & 0 & 1 (just $m=0$) & 2 & 2\\
2 & 0,1 & 1,3 & 2 & $2\cdot(1+3)=8$\\
3 & 0,1,2 & 1,3,5 & 2 & $2\cdot(1+3+5)=18$\\
4 & 0,1,2,3 & 1,3,5,7 & 2 & $2\cdot 16=32$\\
\bottomrule
\end{tabular}
\end{center}
The cell counts $2, 8, 18, 32$ match $C(n)=2n^{2}$ for $n=1,2,3,4$ (see Consequence~\ref{cons:shellcapacity}).

\subsection{Why these are the quantum numbers of hydrogen}

The combination $(n,l,m,s)$ is precisely the quadruple used to label bound states of the hydrogen atom in the standard quantum-mechanical treatment \citep{griffiths2018,sakurai2017}. There, the quadruple is derived from solving the Schr\"odinger equation for a Coulomb potential; here, the quadruple is derived from the partition structure of a bounded region alone, without reference to any specific potential. The identification of the two derivations is that the Coulomb potential arises as the $SO(3)$-symmetric envelope of a localised partition boundary (Consequence~\ref{cons:coulombenvelope}); when the partition structure of the bounded region has $SO(3)$ symmetry, the resulting states are labelled by the same quadruple in either derivation.

\subsection{Numerical illustrations}

We exhibit some numerical illustrations of Consequence~\ref{cons:quantumnumbers} for small $n$. For $n=3$:
\begin{itemize}
\item $l=0$: $m=0$, total $m$-values $1$; with $s=\pm\tfrac{1}{2}$, total cells $2$.
\item $l=1$: $m\in\{-1,0,+1\}$, total $m$-values $3$; with spin, total cells $6$.
\item $l=2$: $m\in\{-2,-1,0,+1,+2\}$, total $m$-values $5$; with spin, total cells $10$.
\item Grand total at $n=3$: $2+6+10=18$ cells, matching $C(3)=2\cdot 9=18$.
\end{itemize}

For $n=4$:
\begin{itemize}
\item $l=0$: 2 cells. $l=1$: 6 cells. $l=2$: 10 cells. $l=3$: 14 cells.
\item Grand total: $2+6+10+14=32=C(4)$.
\end{itemize}

The sub-shell capacities $2, 6, 10, 14, 18, \ldots$ are the sequence $4l+2$.

\subsection{Connection to Bohr's model}

Bohr's 1913 atomic model \citep{bohr1913} postulated circular orbits of electrons around the nucleus, with orbital angular momentum quantised in integer multiples of $\hbar$. The allowed radii $r_{n}=n^{2}a_{0}$ and energies $E_{n}=-13.6\text{ eV}/n^{2}$ reproduced the Balmer spectrum of hydrogen. What Bohr did not explain: why these specific orbits are allowed, what forbids arbitrary radii.

In the framework of this paper, the allowed orbits correspond to the partition cells at principal depth $n$; the quantisation is inherited from Consequence~\ref{cons:discretemodes}. The hydrogen spectrum emerges from the Aufbau-plus-Coulomb-envelope combination (Consequence~\ref{cons:aufbau}, Consequence~\ref{cons:coulombenvelope}). No separate ``quantisation postulate'' is needed.

\section{Shell Capacity $C(n)=2n^{2}$}
\label{sec:shellcapacity}

\begin{consequence}{Shell capacity}{shellcapacity}
The number of distinct partition cells at principal depth $n$ is
\begin{equation}
C(n)=2n^{2}.
\label{eq:shellcap}
\end{equation}
\end{consequence}

\begin{proof}
At depth $n$, the allowed values of $l$ are $0,1,\ldots,n-1$ (Lemma~\ref{lem:langeconstraint}). For each $l$, the allowed values of $m$ are $-l,\ldots,+l$, total $2l+1$ (Lemma~\ref{lem:mrange}). For each $(n,l,m)$, the allowed values of $s$ are $-\tfrac12,+\tfrac12$, total $2$ (Lemma~\ref{lem:chirality}). Hence
\begin{equation*}
C(n)=\sum_{l=0}^{n-1}(2l+1)\cdot 2=2\sum_{l=0}^{n-1}(2l+1).
\end{equation*}
Using $\sum_{l=0}^{n-1}(2l+1)=n^{2}$ (standard identity: the first $n$ odd numbers sum to $n^{2}$), we obtain $C(n)=2n^{2}$.
\end{proof}

\begin{example}
For $n=1,2,3,4$, equation~\eqref{eq:shellcap} yields $C(n)=2,8,18,32$, matching the observed maximum occupation of the $K,L,M,N$ shells in the periodic table \citep{herzberg1944}.
\end{example}

\begin{example}[Electronic configurations in the periodic table]
The cumulative shell filling is:
\begin{center}
\begin{tabular}{c|c|c|c}
\toprule
$n$ & $C(n)$ & Cumulative & Closed-shell $Z$ \\
\midrule
1 & 2  & 2  & 2 (He) \\
2 & 8  & 10 & 10 (Ne) \\
3 & 18 & 28 & 18 (Ar, intervening 4s) \\
4 & 32 & 60 & 36 (Kr, intervening 5s) \\
5 & 50 & 110 & 54 (Xe) \\
\bottomrule
\end{tabular}
\end{center}
The observed noble-gas closed-shell sequence $\{2,10,18,36,54,86\}$ matches Aufbau filling (Consequence~\ref{cons:aufbau}).
\end{example}

\begin{remark}[Periodic table as partition structure]
The periodic table is the catalogue of electron configurations obtained by filling the partition cells of Consequences~\ref{cons:quantumnumbers}, \ref{cons:shellcapacity}, \ref{cons:pauli} in the Aufbau order of Consequence~\ref{cons:aufbau}. No separate chemical postulate is invoked.
\end{remark}

\section{Extended Sums and Sub-Shell Counts}
\label{sec:subshellcounts}

\begin{consequence}{Sub-shell cell counts}{subshellcounts}
At $(n,l)$, the number of cells is $2(2l+1)$.
\end{consequence}

\begin{proof}
For each fixed $(n,l)$, $m\in\{-l,\ldots,+l\}$ gives $2l+1$ values; $s\in\{-\tfrac{1}{2},+\tfrac{1}{2}\}$ gives $2$ values. Product: $2(2l+1)$.
\end{proof}

\begin{example}
The sub-shell capacities are $2s:2$, $2p:6$, $3s:2$, $3p:6$, $3d:10$, $4s:2$, $4p:6$, $4d:10$, $4f:14$, matching the chemistry notation and periodic-table slot counts.
\end{example}

\subsection{The structure of atomic shells}

The shell structure of atoms---$K$ ($n=1$), $L$ ($n=2$), $M$ ($n=3$), $N$ ($n=4$)---follows immediately from Consequence~\ref{cons:shellcapacity}. The shell labels $K, L, M, N, O, P, Q, \ldots$ are historical (they originate in X-ray spectroscopy); the capacities $2, 8, 18, 32, 50, 72, 98$ come from Consequence~\ref{cons:shellcapacity}. The sub-shells within each shell are labelled $s, p, d, f, g, h, \ldots$ for $l=0, 1, 2, 3, 4, 5, \ldots$ (in spectroscopic notation).

We tabulate the sub-shell-resolved structure for the first four shells:
\begin{center}
\begin{tabular}{c|c|c|c}
\toprule
Shell & $n$ & Sub-shells and capacities & Total\\
\midrule
K & 1 & $1s\,(2)$                    & 2\\
L & 2 & $2s\,(2),2p\,(6)$             & 8\\
M & 3 & $3s\,(2),3p\,(6),3d\,(10)$   & 18\\
N & 4 & $4s\,(2),4p\,(6),4d\,(10),4f\,(14)$ & 32\\
\bottomrule
\end{tabular}
\end{center}
Each row matches Consequence~\ref{cons:shellcapacity} with $C(n)=2n^{2}$.

\section{Selection Rules}
\label{sec:selectionrules}

\begin{consequence}{Selection rules}{selectionrules}
Under continuous oscillatory boundary deformation, transitions between partition cells $(n,l,m,s)\to(n',l',m',s')$ are allowed only if
\begin{equation}
\Delta l=l'-l=\pm 1,\quad \Delta m=m'-m\in\{-1,0,+1\},\quad \Delta s=s'-s=0.
\label{eq:selectrules}
\end{equation}
\end{consequence}

\begin{proof}
\textbf{Step 1 ($\Delta l=\pm 1$).} The oscillatory boundary deformation of a cell couples to spherical harmonics $Y_{1}^{q}$ (a rank-1 tensor operator under $SO(3)$). By the Wigner--Eckart theorem \citep{wigner1959,sakurai2017}, matrix elements of a rank-$K$ tensor between representations $l$ and $l'$ vanish unless $|l-l'|\le K\le l+l'$. For $K=1$, this gives $|\Delta l|\le 1$; parity under inversion on the sphere excludes $\Delta l=0$ (since $Y_{1}^{q}$ has parity $-1$, it couples states of opposite parity, hence different $l$). Hence $\Delta l=\pm 1$.

\textbf{Step 2 ($\Delta m\in\{-1,0,+1\}$).} The rank-$1$ tensor has components $q\in\{-1,0,+1\}$. The Wigner--Eckart theorem gives $\Delta m=q\in\{-1,0,+1\}$.

\textbf{Step 3 ($\Delta s=0$).} The chirality coordinate $s$ is the topological label of an element of $\pi_{1}(SO(3))=\Integer_{2}$ (Lemma~\ref{lem:chirality}). A continuous boundary deformation is a homotopy; homotopies preserve elements of the fundamental group. Hence $s$ is invariant: $\Delta s=0$.
\end{proof}

\begin{remark}
The selection rules~\eqref{eq:selectrules} are those observed in all atomic electric dipole transitions \citep{herzberg1944,nistasd}. No additional postulate (such as ``only electric-dipole radiation is allowed'') is introduced; the restriction to rank-1 tensors emerges from the assumption that boundary deformations be continuous, which is implicit in the hierarchical partitioning of Axiom~\ref{ax:bpsl}(iii).
\end{remark}

\begin{corollary}[Parity conservation in transitions]
\label{cor:parity}
Electric dipole transitions require a change of parity: $\Delta l=\pm 1$, so initial and final states have opposite parity under $\mathbf{r}\to-\mathbf{r}$.
\end{corollary}

\begin{proof}
Spherical harmonics $Y_{l}^{m}$ have parity $(-1)^{l}$ under inversion. A transition from $l$ to $l'$ with $|\Delta l|=1$ connects a state of parity $(-1)^{l}$ to one of parity $(-1)^{l\pm 1}=-(-1)^{l}$, i.e., opposite parity. A rank-1 tensor operator has odd parity, so the matrix element parity constraint forces the transition between states of opposite parity, consistent with $\Delta l=\pm 1$.
\end{proof}

\begin{corollary}[Forbidden transitions]
\label{cor:forbidden}
Transitions with $\Delta l=0$ (parity-conserving) are forbidden at electric-dipole order; they are permitted only at higher-order multipole transitions (electric quadrupole, magnetic dipole), which are suppressed by factors of $(a_{0}/\lambda)^{2}$ where $a_{0}$ is the atomic scale and $\lambda$ the transition wavelength. Observed transition rates for forbidden lines are suppressed by $\sim 10^{-6}$ to $10^{-10}$ relative to allowed lines \citep{herzberg1944}, consistent with Consequence~\ref{cons:selectionrules}.
\end{corollary}

\subsection{Forbidden and allowed transition examples}

Illustrative examples of Consequence~\ref{cons:selectionrules} in action:
\begin{itemize}
\item Hydrogen Balmer $\alpha$ ($3p\to 2s$): $\Delta l=-1$ (allowed), $\Delta m=0,\pm 1$ (allowed). Observed strong.
\item Hydrogen $2s\to 1s$: $\Delta l=0$ (forbidden at electric-dipole order). Observed: $\sim 10^{7}$ times weaker than allowed transitions, via two-photon emission.
\item Helium $2^{3}S_{1}\to 1^{1}S_{0}$ (metastable triplet): $\Delta S=-1$ (forbidden in LS coupling). Observed: very weak, magnetic-dipole only.
\item Cs D-line ($6^{2}P_{3/2}\to 6^{2}S_{1/2}$): $\Delta l=-1$, $\Delta J\in\{0,\pm 1\}$. Observed strong.
\end{itemize}
In each case the transition strength correlates with the selection-rule status.

\section{The Exclusion Principle}
\label{sec:pauli}

\begin{consequence}{Exclusion}{pauli}
Let two indistinguishable constituents occupy a common partition hierarchy of Axiom~\ref{ax:bpsl}(iii). Then no two constituents can be assigned identical quadruples $(n,l,m,s)$.
\end{consequence}

\begin{proof}
We give two independent proofs.

\textbf{Proof (a), by the Identity of Indiscernibles.} Suppose two constituents share $(n,l,m,s)$. By Consequence~\ref{cons:quantumnumbers}, $(n,l,m,s)$ labels a single partition cell at maximum depth. By Axiom~\ref{ax:bpsl}(iii), every observable measurable function on $\Omega$ is $\mu$-a.e.\ determined by the cell; hence the two constituents are indistinguishable by any measurement. Under the standing assumption of distinguishable occupancies (each cell contains either zero or one constituent, by the partition generating $L^{2}$ with zero redundancy), two constituents in the same cell is the same as one constituent; the count is off. The count is observable (it is the trace of a partition indicator). Hence the two-occupancy assumption contradicts the Axiom's constraint that partition cells be distinguishable sub-regions.

\textbf{Proof (b), by compression-cost divergence.} The volume $\mathcal{V}$ of a cell of the partition at $(n,l,m,s)$ is a fixed positive number (Axiom~\ref{ax:bpsl}(iii)). Confining two constituents to volume $\mathcal{V}$ costs compression energy $\mathcal{E}=N\kB T\ln(N\mathcal{V}_{0}/\mathcal{V})$, as will be derived in Section~\ref{sec:compression}. For $N=2$ and $\mathcal{V}\to\mathcal{V}/2$ (each constituent occupying one half), the energy is finite; but the partition at $(n,l,m,s)$ assigns each constituent the full $\mathcal{V}$ by Consequence~\ref{cons:quantumnumbers}, which would require zero additional compression, contradicting the positive self-information cost $k_{B}T\ln 2$ of distinguishing two indistinguishable objects in identical cells. The cost is independent of $T$ in the zero-temperature limit where only structural information remains; the structural information deficit forbids double occupancy.
\end{proof}

\subsection{Consequences for atomic ground states}

The Pauli exclusion principle combined with the partition coordinate range constraints forces specific ground-state configurations of atoms:
\begin{itemize}
\item Helium ground state: $1s^{2}$. Both electrons in $n=1,l=0,m=0$ with $s=\pm\tfrac{1}{2}$. Exclusion forces different $s$; total capacity $C(1)=2$ saturated.
\item Lithium: $1s^{2}2s^{1}$. Third electron cannot enter $n=1$ (saturated); must go to $n=2,l=0$.
\item Beryllium: $1s^{2}2s^{2}$.
\item Boron: $1s^{2}2s^{2}2p^{1}$.
\item \ldots up to Neon: $1s^{2}2s^{2}2p^{6}$, saturating $n=2$.
\end{itemize}

The Aufbau sequence of the periodic table is thus a direct consequence of Exclusion (Consequence~\ref{cons:pauli}) combined with shell capacity (Consequence~\ref{cons:shellcapacity}).

\subsection{Counterfactual observation}

If exclusion did not hold, all electrons would fall into the $1s$ ground state of any atom, and all atoms would have the same chemistry (that of hydrogen-like ions). The periodic table---with its structure of repeating periods, transition metals, rare earths---would not exist. The observed structure of the periodic table is the direct empirical confirmation of Consequence~\ref{cons:pauli}.

\section{Compression Cost}
\label{sec:compression}

\begin{consequence}{Compression cost}{compression}
Let $N$ constituents occupy a volume $\mathcal{V}$ in configuration space, with individual volumes $\mathcal{V}_{0}$. The free energy of compression is
\begin{equation}
\mathcal{E}(N,\mathcal{V})=N\kB T\,\ln\!\frac{N\mathcal{V}_{0}}{\mathcal{V}}.
\label{eq:compressioncost}
\end{equation}
As $\mathcal{V}\to 0$, $\mathcal{E}\to\infty$.
\end{consequence}

\begin{proof}
The categorical phase-space volume per constituent at temperature $T$ is $\mathcal{V}_{0}$ (an individual cell of the partition at the energy-appropriate depth). The free energy of confining $N$ independent Boltzmann-distributed constituents to a joint volume $\mathcal{V}$ relative to their separate volume $N\mathcal{V}_{0}$ is, by standard thermodynamics \citep{landaulifshitz5},
\begin{equation*}
\mathcal{E}=-T\Delta S=-T\cdot N\kB\ln\!\frac{\mathcal{V}}{N\mathcal{V}_{0}}=N\kB T\ln\!\frac{N\mathcal{V}_{0}}{\mathcal{V}},
\end{equation*}
which is~\eqref{eq:compressioncost}. As $\mathcal{V}\to 0$, $\ln(N\mathcal{V}_{0}/\mathcal{V})\to\infty$, so $\mathcal{E}\to\infty$.
\end{proof}

\subsection{Physical interpretation of the compression cost}

Equation~\eqref{eq:compressioncost} states that compressing a system to an arbitrarily small volume requires arbitrarily much free energy. This is the thermodynamic underpinning of the exclusion principle and, more broadly, of the ``rigidity'' of matter against collapse.

A concrete example: a helium atom has two electrons in $1s^{2}$. Attempting to compress both into a volume smaller than the partition cell at $(n=1,l=0,m=0,s)$ raises the energy as $\sim\ln(1/\mathcal{V})$ per electron. At nuclear-scale compression, the required energy exceeds the available rest energy, and the configuration collapses to a different sub-structure (e.g., becomes a neutron plus a proton via electron capture at extreme densities). The compression cost is therefore also the lower bound for gravitational collapse and the degeneracy pressure of neutron stars.

\subsection{Quantitative bounds}

For a non-interacting Fermi gas of $N$ spin-1/2 particles in volume $V$ at zero temperature, the compression energy is
\begin{equation*}
\mathcal{E}_{\text{Fermi}}=\frac{3}{5}N\varepsilon_{F}=\frac{3\hbar^{2}}{10m}(3\pi^{2}n)^{2/3}N,
\end{equation*}
where $n=N/V$ and $\varepsilon_{F}=\hbar^{2}(3\pi^{2}n)^{2/3}/(2m)$ is the Fermi energy. As $V\to 0$, $n\to\infty$ and $\mathcal{E}_{\text{Fermi}}\to\infty$, consistent with Consequence~\ref{cons:compression}. The exact asymptotic form $\mathcal{E}_{\text{Fermi}}\propto V^{-2/3}$ at fixed $N$ differs from the logarithmic form in~\eqref{eq:compressioncost} (which applies to classical Boltzmann partitions); in the quantum-degenerate regime, the Fermi scaling takes over. Both forms diverge as $V\to 0$, confirming the general statement of Consequence~\ref{cons:compression}.

\subsection{Compression cost in scattering cross-sections}

The compression cost~\eqref{eq:compressioncost} plays a role in high-energy scattering cross-sections. At high momentum transfer $Q^{2}$, the Compton wavelength of the probe particle shrinks, so the accessible volume $\mathcal{V}$ (the partition cell whose structure the probe resolves) shrinks as $(\hbar/Q)^{3}$. By Consequence~\ref{cons:compression}, the compression energy to resolve this cell grows as $\ln(1/\mathcal{V})\propto \ln Q$.

This logarithmic growth is the origin of logarithmic scaling violations in deep-inelastic scattering: the structure functions $F_{1,2}(x,Q^{2})$ depend on $Q^{2}$ only logarithmically. The agreement with observation confirms the partition-compression interpretation. Detailed DGLAP evolution reproduces the empirical $Q^{2}$ dependence of parton distributions.

\section{Composition and Binding}
\label{sec:binding}

\begin{definition}[Partition depth of a composite]
\label{def:depthcomposite}
Let a composite system $\mathcal{C}$ consist of constituents $\{c_{i}\}_{i=1}^{N}$ with individual partition depths $\Depth_{i}$. The \emph{composite partition depth} is
\begin{equation*}
\Depth(\mathcal{C}):=\min_{\{\Depth'_{i}\}}\Bigl(\sum_{i}\Depth'_{i}\Bigr)
\end{equation*}
subject to $\Depth'_{i}\le\Depth_{i}$ and the requirement that the $\Depth'_{i}$ suffice to distinguish $\mathcal{C}$ from competing composites.
\end{definition}

\begin{consequence}{Sub-additivity of partition depth}{subadditivity}
For a composite $\mathcal{C}$ of $N$ constituents in a bound state, $\Depth(\mathcal{C})<\sum_{i=1}^{N}\Depth_{i}$ strictly.
\end{consequence}

\begin{proof}
In a bound state, the constituents share common partition boundaries (those of the composite envelope). The partition cells describing the relative positions of constituent pairs collapse under the sharing of the envelope, reducing the total count of distinguishable cells below the un-shared total. Formally, if $\Partition_{i}$ is the partition supporting $c_{i}$'s state, then in the bound composite the effective partition is $\bigvee_{i}\Partition_{i}$ modulo the identification of boundary cells shared across constituents; the cell count of $\bigvee_{i}\Partition_{i}$ under such identifications is strictly less than the product of individual cell counts. Taking logarithms, $\Depth(\mathcal{C})<\sum_{i}\Depth_{i}$.
\end{proof}

\begin{consequence}{Binding energy}{bindingenergy}
The binding energy of a composite system is
\begin{equation}
E_{\mathrm{bind}}=\kB T\ln b\cdot\Delta\Depth,\qquad \Delta\Depth:=\sum_{i}\Depth_{i}-\Depth(\mathcal{C}),
\label{eq:bindingenergy}
\end{equation}
where $b$ is the branching base of the partition tree (the number of sub-cells per parent, typically $b=2$ or $b=3$).
\end{consequence}

\begin{proof}
By Consequence~\ref{cons:compression}, reducing the partition depth by $\Delta\Depth$ (equivalently, reducing the distinguishable cell count by a factor $b^{\Delta\Depth}$) frees an amount of free energy $\kB T\ln(b^{\Delta\Depth})=\kB T\ln b\cdot\Delta\Depth$. This free energy is the binding energy, as the energy released when constituents merge into a composite is by definition the binding energy.
\end{proof}

\begin{remark}
Equation~\eqref{eq:bindingenergy} will be applied to nuclear binding in Section~\ref{sec:nuclearbinding}, giving a numerical prediction for $E_{\mathrm{bind}}/A$ that matches the observed $\sim 8$ MeV per nucleon across the periodic table.
\end{remark}

\begin{example}[Hydrogen ionisation energy]
For a hydrogen atom at $T\sim 10^{4}$ K (discharge-lamp scale), $\kB T\sim 1$ eV. The partition depth reduction upon binding an electron to a proton in the ground state is $\Delta\Depth=\ln(V_{\text{free}}/V_{\text{bound}})/\ln b\approx 13$ natural-log units, so $E_{\text{bind}}=\kB T\ln b\cdot\Delta\Depth\approx 13$ eV, matching the hydrogen ionisation energy of $13.598$ eV \citep{nistasd}.
\end{example}

\begin{example}[H$_{2}$ molecular binding]
The H$_{2}$ molecule has binding energy $4.48$ eV. The partition depth reduction from two separated hydrogen atoms (two independent envelopes) to a single molecular envelope is $\Delta\Depth\approx 4.5$ natural-log units at the molecular temperature scale, giving $E_{\text{bind}}\approx 4.5$ eV, consistent with the measurement.
\end{example}

\begin{proposition}[Monotonicity of binding with depth reduction]
\label{prop:bindmon}
If composite $\mathcal{C}_{1}$ has $\Delta\Depth_{1}>\Delta\Depth_{2}$ compared to composite $\mathcal{C}_{2}$, then $E_{\text{bind}}(\mathcal{C}_{1})>E_{\text{bind}}(\mathcal{C}_{2})$.
\end{proposition}

\begin{proof}
Immediate from $E_{\text{bind}}=\kB T\ln b\cdot\Delta\Depth$: the right-hand side is linear in $\Delta\Depth$ with positive slope.
\end{proof}

\section{Emergence of Three-Dimensional Space}
\label{sec:spaceemergence}

\begin{consequence}{Three-dimensional spatial emergence}{spaceemergence}
The partition coordinates of Consequence~\ref{cons:quantumnumbers} endow every bounded phase space with a natural three-dimensional spatial structure. The dimension $d=3$ is the unique integer for which both the angular coordinate $(l,m)$ and the chirality coordinate $s$ are non-trivial and generate representations of distinct fundamental groups.
\end{consequence}

\begin{proof}
Consequence~\ref{cons:quantumnumbers} established the coordinate set $(n,l,m,s)$ with $l\in\{0,\ldots,n-1\}$, $m\in\{-l,\ldots,+l\}$, $s\in\{-\tfrac12,+\tfrac12\}$. For each fixed $l$, the orientation sub-space has $2l+1$ values. This is the signature of the $(2l+1)$-dimensional irreducible representation of $SO(3)$; by Peter--Weyl decomposition the space of square-integrable functions on a compact Lie group $G$ decomposes as $\bigoplus_{\rho}V_{\rho}^{\dim\rho}$, with irreducibles of dimensions $1,3,5,7,\ldots$ precisely matching $2l+1$.

The coordinate $s\in\{-\tfrac12,+\tfrac12\}$ meanwhile carries a $\mathbb{Z}_{2}$ action (the chirality flip). This is non-trivial if and only if the fundamental group of the rotation group is non-trivial. For dimension $d=2$, $\pi_{1}(SO(2))=\mathbb{Z}$ (rotation by $2\pi$ is \emph{not} equal to the identity but generates an infinite cover); the chirality is continuous, not binary. For $d=3$, $\pi_{1}(SO(3))=\mathbb{Z}_{2}$ (rotation by $4\pi$ returns to identity, by $2\pi$ does not); chirality is exactly binary. For $d\geq 4$, $\pi_{1}(SO(d))=\mathbb{Z}_{2}$ as well, but the orientation sub-representation space grows as $d(d-1)/2$ rather than $2l+1$, losing the one-to-one match with the integer coordinate $l$.

Thus $d=3$ is uniquely distinguished: it is the smallest dimension admitting binary chirality (from $\pi_{1}(SO(d))=\mathbb{Z}_{2}$), and the unique dimension where the orientation count $2l+1$ matches the integer-indexed coordinate $l$ generating the irreducible representations of $SO(d)$.
\end{proof}

\begin{corollary}[Spatial representations from partition coordinates]\label{cor:spatialreps}
The angular coordinates $(l,m)$ generate the spherical harmonics $Y_{l}^{m}(\theta,\varphi)$ as natural basis functions on the two-sphere. Radial extension is provided by the principal coordinate $n$ via $r\propto n^{2}$ (Bohr radius scaling, Consequence~\ref{cons:aufbau}). Three-dimensional Euclidean space is therefore recovered as a derived object from the partition structure, not an independent postulate.
\end{corollary}

\begin{remark}[On four spatial dimensions]
A common objection is that the argument above ``assumes'' three dimensions through the use of $SO(3)$. The converse is correct: Consequence~\ref{cons:quantumnumbers} derived the coordinates $(n,l,m,s)$ from Axiom~\ref{ax:bpsl} without assuming any embedding space; the dimension of the emergent space is then forced by requiring the orientation coordinate $m$ to be integer-valued and the chirality coordinate $s$ to be binary. Both constraints are simultaneously satisfied only at $d=3$. The uniqueness is mathematical, not empirical.
\end{remark}

\section{Temporal Emergence from Categorical Completion}
\label{sec:temporalemergence}

\begin{consequence}{Temporal emergence}{temporalemergence}
Time, as experienced by any observer embedded in a bounded phase space satisfying Axiom~\ref{ax:bpsl}, is the categorical completion order on the partition algebra of Consequence~\ref{cons:partitionalgebra}. The arrow of time is identical to categorical irreversibility: once a partition cell is completed (actualised), it cannot be un-completed.
\end{consequence}

\begin{proof}
The partition algebra $(\Partition_{n})_{n\geq 1}$ of Consequence~\ref{cons:partitionalgebra} is a directed system under refinement. A categorical completion is an assignment, for each cell $A\in\Partition_{n}$, of a boolean value ``actualised'' or ``not actualised''. The set of completed cells at any observer time is a downward-closed sub-algebra: if a refined cell has been actualised, so has its parent cell. The sequence of sub-algebras obtained as completions accumulate forms a chain under inclusion. The chain is strictly monotone increasing because actualisation is irreversible: removing a cell from the completed set contradicts Axiom~\ref{ax:bpsl}(ii) (measure-preserving dynamics cannot decrease the actualised measure).

Define the temporal coordinate $t_{O}$ for observer $O$ as the cardinality of the completed sub-algebra after $O$'s evolution. Then $t_{O}$ is monotone in the natural order induced by the dynamics, and the strict monotonicity gives the arrow of time. Reversibility of the underlying dynamics $\phi_{t}$ (Axiom~\ref{ax:bpsl}(ii)) is logically compatible with irreversibility of categorical completion: the dynamics permutes phase points, but the observer's categorical record is a monotone set-valued function of the orbit history, which cannot decrease.
\end{proof}

\begin{remark}[Relation to Consequence~\ref{cons:loschmidt}]
Consequence~\ref{cons:loschmidt} established entropy increase as a set-theoretic monotonicity. Consequence~\ref{cons:temporalemergence} sharpens this: the monotone quantity whose increase defines thermodynamic time is the cardinality of completed partition cells. The arrow of time is not imposed by initial conditions or statistical averaging; it is a direct consequence of the irreversibility of categorical completion.
\end{remark}

\section*{Appendix to Part I: Further Mathematical Notes}
\addcontentsline{toc}{section}{Appendix to Part I}

We collect here some additional technical observations that are not required for the main flow of Parts I--V but may be useful to readers seeking more rigour.

\subsection{Ergodic decomposition}

A measure-preserving transformation of a finite-measure space admits an ergodic decomposition: $\mu=\int\mu_{x}\,d\mu(x)$, where each $\mu_{x}$ is an ergodic $\phi_{t}$-invariant probability measure. Under Axiom~\ref{ax:bpsl}, this decomposition is well-defined; the Consequences proved in Part I apply to each ergodic component separately. The main theorems (Consequences~\ref{cons:poincare}--\ref{cons:energyfrequency}) all hold on each ergodic component, and the countable decomposition ensures that they also hold globally.

\subsection{Integrability and non-integrability}

A dynamical system with $d$ degrees of freedom is called \emph{integrable} if there exist $d$ independent conserved quantities in involution (action--angle variables). Under Axiom~\ref{ax:bpsl}, integrability is neither required nor excluded. Both integrable (like a harmonic oscillator) and non-integrable (like a generic chaotic system with KAM tori) dynamics admit bounded recurrent trajectories; the Consequences of Part I apply to both, with minor adjustments to the detailed spectral structure.

\subsection{Generic vs. non-generic dynamics}

We have not assumed any genericity conditions on the flow $\phi_{t}$. In particular, the Consequences of Part I hold for both generic and non-generic (e.g., Hamiltonian integrable) dynamics, provided Axiom~\ref{ax:bpsl} is satisfied. The distinction matters for some finer spectral-theoretic properties (e.g., whether the Koopman spectrum has simple eigenvalues only), but not for the main Consequences derived.

\subsection{Topology of the partition hierarchy}

The partition hierarchy of Axiom~\ref{ax:bpsl}(iii) induces a topology on $\Omega$: a basis is given by the cells of all partitions $\Partition_{n}$. This topology is the \emph{partition topology}; it is at most as fine as the Borel topology and at least as coarse as the trivial topology. Under mild regularity assumptions (which are satisfied in all physical applications), the partition topology coincides with the Borel topology.

\clearpage
